% Prezentare licenta ManFast - Beamer
% Compilare: pdflatex prezentare.tex (de 2 ori)
\documentclass[aspectratio=169,12pt]{beamer}

\usepackage[utf8]{inputenc}
\usepackage[T1]{fontenc}
\usepackage[romanian]{babel}
\usepackage{graphicx}
\usepackage{booktabs}
\usepackage{hyperref}

\usetheme{Madrid}
\usecolortheme{whale}
\setbeamertemplate{navigation symbols}{}
\setbeamertemplate{footline}[frame number]

\graphicspath{{figuri/}{../diagrame/}}

% Incarca imagine daca exista; altfel afiseaza placeholder text
\newcommand{\slideimage}[3]{%
	\begin{center}
		\IfFileExists{#1}{%
			\includegraphics[width=#2]{#1}%
		}{%
			\fbox{\parbox[c][#3][c]{0.85\textwidth}{%
					\centering\small\textit{[Inserează poza]}\\[0.2em]
					\small\texttt{\detokenize{#1}}%
			}}%
		}%
	\end{center}%
}

\title{Aplicație pentru organizarea și gestionarea evenimentelor}
\subtitle{Implementare: ManFast}
\author{Ene Luis Adrian}
\institute{%
	Facultatea de Mathematică și Informatică\\
	Specializarea Informatică\\
	Universitatea Ovidius din Constanța
}
\date{2026}

\begin{document}
	
	% --- Slide 1: Titlu ---
	\begin{frame}
		\titlepage
		\vspace{-0.5cm}
		\begin{center}
			\small Coordonator științific: Lect. univ. dr. Rusu Andrei
		\end{center}
	\end{frame}
	
	% --- Slide 2: Structura ---
	\begin{frame}{Structura prezentării}
		\begin{itemize}
			\item Context și obiective
			\item Soluția propusă (ManFast)
			\item Arhitectură și tehnologii
			\item Funcționalități utilizator și administrator
			\item Flux principal și model de date
			\item Rezultate obținute
			\item Demonstrație live
			\item Concluzii și perspective
		\end{itemize}
	\end{frame}
	
	% --- Slide 3: Problema ---
	\begin{frame}{Context și problematica actuală}
		\begin{block}{Ideea principală}
			Organizarea evenimentelor pe piața locală rămâne adesea fragmentată și dependentă de procese manuale.
		\end{block}
		\vspace{0.3cm}
		\begin{itemize}
			\item Promovare dispersată pe canale media diferite fără o centralizare clară.
			\item Plăți nesigure sau manuale (transferuri bancare directe, cash la locație).
			\item Gestionarea participanților pe liste fizice (hârtie), generând cozi la intrare.
			\item Dificultăți în reconcilierea numărului de locuri disponibile cu plățile confirmate.
			\item Lipsa unei platforme dedicate adaptate complet pieței naționale (monedă locală, județe).
		\end{itemize}
	\end{frame}
	
	% --- Slide 4: Obiective ---
	\begin{frame}{Obiectivele lucrării}
		\begin{enumerate}
			\item \textbf{Securitate:} Autentificare robustă a utilizatorilor (validare email, token JWT, resetare parolă).
			\item \textbf{Experiență Fluidă:} Descoperirea rapidă a evenimentelor prin mecanisme de filtrare avansată.
			\item \textbf{Gestiune Financiară:} Emitere de bilete nominale și integrarea plăților online securizate prin Stripe (în RON).
			\item \textbf{Control Centralizat:} Panou administrativ dedicat pentru managementul evenimentelor și validarea rapidă a participanților (check-in).
			\item \textbf{Rigurozitate Tehnică:} Documentarea completă a sistemului prin diagrame de arhitectură, diagrame UML și modele ERD.
		\end{enumerate}
	\end{frame}
	
	% --- Slide 5: Solutia ---
	\begin{frame}{Soluția propusă: Platforma ManFast}
		\begin{block}{Platformă web completă full-stack}
			Un ecosistem digital unificat, dedicat organizatorilor și participanților la conferințe, meeting-uri, spectacole sau concerte.
		\end{block}
		\vspace{0.4cm}
		\begin{itemize}
			\item \textbf{Evenimente centralizate:} Toate detaliile organizatorice sunt disponibile într-un singur punct de acces.
			\item \textbf{Bilete nominale:} Fiecare tichet este unic asociat unui cont securizat, prevenind fraudele.
			\item \textbf{Plăți online rapide:} Integrare nativă cu Stripe Checkout pentru o experiență de plată modernă.
			\item \textbf{Check-in digital:} Interfață administrativă optimizată pentru scanarea și validarea biletelor la intrare.
		\end{itemize}
	\end{frame}
	
	% --- Slide 6: Arhitectura (Imagine Pastrata) ---
	\begin{frame}{Arhitectura sistemului}
		\begin{center}
			\textbf{Model Client-Server cu Organizare Stratificată}
		\end{center}
		\vspace{-0.2cm}
		% POZA PASTRATA CONFORM CERINTEI
		\slideimage{Arhitectura sistemului ManFast.jpg}{0.675\textwidth}{3.3cm}
		\vspace{-0.2cm}
		\begin{center}
			\small React Frontend \textrightarrow{} Express API Backend \textrightarrow{} PostgreSQL Database \quad|\quad Servicii externe: Stripe, Brevo
		\end{center}
	\end{frame}
	
	% --- Slide 7: Tehnologii ---
	\begin{frame}{Tehnologii utilizate}
		\begin{block}{Stiva tehnologică completă}
			\centering
			\vspace{0.1cm}
			\begin{tabular}{@{}ll@{}}
				\toprule
				\textbf{Strat / Componentă} & \textbf{Tehnologii alese} \\
				\midrule
				Frontend & React 19, Vite, Material UI 9 \\
				Backend & Node.js, Express API, JWT, bcrypt \\
				Bază de date & PostgreSQL \\
				Procesare Plăți & Stripe Checkout API (Procesare nativă în RON) \\
				Serviciu Email & Brevo SMTP (Verificare cont și notificări) \\
				\bottomrule
			\end{tabular}
		\end{block}
	\end{frame}
	
	% --- Slide 8: Utilizator ---
	\begin{frame}{Funcționalități -- Rolul de Utilizator}
		\begin{columns}[T]
			\column{0.48\textwidth}
			\begin{block}{Acces și Securitate}
				\begin{itemize}
					\item Înregistrare cont nou.
					\item Verificare obligatorie prin cod e-mail.
					\item Autentificare securizată prin token-uri JWT stocate rigid.
					\item Gestiunea profilului și resetarea parolei.
				\end{itemize}
			\end{block}
			\column{0.48\textwidth}
			\begin{block}{Explorare și Achiziție}
				\begin{itemize}
					\item Căutare și filtrare avansată a evenimentelor după multiple criterii.
					\item Achiziție securizată de bilete.
					\item Zonă dedicată: „Biletele mele”.
					\item Generare automată a datelor nominale de acces.
				\end{itemize}
			\end{block}
		\end{columns}
	\end{frame}
	
	% --- Slide 9: Administrator ---
	\begin{frame}{Funcționalități -- Rolul de Administrator}
		\begin{columns}[T]
			\column{0.48\textwidth}
			\begin{block}{Management de Conținut}
				\begin{itemize}
					\item Creare completă de evenimente noi.
					\item Încărcare de imagini reprezentative.
					\item Configurare detalii logistice, dată, timp și preț bilet.
					\item Editarea și actualizarea datelor.
				\end{itemize}
			\end{block}
			\column{0.48\textwidth}
			\begin{block}{Control și Validare}
				\begin{itemize}
					\item Vizualizarea listei complete de participanți înscriși.
					\item Modul dedicat de check-in în timp real la intrarea în sală.
					\item Drepturi avansate de promovare sau revocare a altor administratori.
				\end{itemize}
			\end{block}
		\end{columns}
	\end{frame}
	
	% --- Slide 10: Flux principal ---
	\begin{frame}{Fluxul principal al unui participant}
		\begin{block}{Traseul utilizatorului în platformă}
			\begin{enumerate}
				\item \textbf{Descoperire:} Identifică și filtrează evenimentul dorit în interfața principală.
				\item \textbf{Autentificare:} Se conectează securizat sau confirmă contul nou prin e-mail.
				\item \textbf{Tranzacționare:} Selectează biletul nominal și inițiază fluxul securizat Stripe Checkout.
				\item \textbf{Posesie:} După confirmarea plății, biletul devine instant disponibil în secțiunea personală.
				\item \textbf{Validare:} La locație, administratorul validează biletul prin sistemul intern de check-in.
			\end{enumerate}
		\end{block}
	\end{frame}
	
	% --- Slide 11: Model date (Imagine Pastrata) ---
	\begin{frame}{Modelul de date (ERD)}
		\vspace{-0.2cm}
		% POZA PASTRATA CONFORM CERINTEI
		\slideimage{Model de domeniu _ ERD ManFast.jpg}{0.75\textwidth}{4cm}
		\vspace{-0.2cm}
		\begin{center}
			\small Structură relațională optimizată. Tabele nucleu: \texttt{users}, \texttt{events}, \texttt{tickets}
		\end{center}
	\end{frame}
	
	% --- Slide 12: Rezultate ---
	\begin{frame}{Rezultate obținute}
		\begin{itemize}
			\item \textbf{Aplicație completă:} Mediu full-stack perfect funcțional, fluid și stabil.
			\item \textbf{Securitate end-to-end:} Sistem de autentificare modern bazat pe cele mai bune practici din industrie.
			\item \textbf{Tranzacții reale:} Sistem de plăți sandbox complet funcțional în moneda națională (RON).
			\item \textbf{Panou administrativ integrat:} Gestiune dinamică și proces rapid de validare la eveniment.
			\item \textbf{Bază solidă documentară:} Arhitectură curată, diagrame UML clare și API REST documentat extensiv.
		\end{itemize}
	\end{frame}
	
	% --- Slide 13: Demonstratie ---
	\begin{frame}{Demonstrație practică}
		\begin{columns}[T]
			\column{0.48\textwidth}
			\begin{block}{Videoclip demonstrativ}
				\begin{itemize}
					\item Videoclip
				\end{itemize}
			\end{block}
		\end{columns}
	\end{frame}
	
	% --- Slide 14: Concluzii ---
	\begin{frame}{Concluzii și perspective de dezvoltare}
		\begin{block}{Concluzii}
			Rezultatul: o platformă modernă, robustă și adaptată cerințelor actuale.
		\end{block}
		\vspace{0.2cm}
		\textbf{Direcții viitoare de dezvoltare (Perspective):}
		\begin{itemize}
			\item Implementarea unui sistem avansat de Webhooks Stripe pentru confirmarea asincronă a plăților.
			\item Suport pentru categorii multiple de bilete (VIP, Early Bird, Acces General).
			\item Integrarea testelor automate unitare și end-to-end
		\end{itemize}
	\end{frame}
	
	% --- Slide 15: Multumiri ---
	\begin{frame}{Încheiere}
		\begin{center}
			\Large \textbf{Vă mulțumesc pentru atenție!}
			\vspace{0.8cm}
			
			\normalsize
			\textbf{Lect. univ. dr. Rusu Andrei}\\
			(Coordonator științific)
			\vspace{0.5cm}
			
			Membrii comisiei de examinare\\
		
		\end{center}
	\end{frame}
	
	% --- Slide 16: Bibliografie ---
	\begin{frame}[allowframebreaks]{Bibliografie selectivă}
		\footnotesize
		\begin{itemize}
			\item R. Fielding. \textit{Architectural Styles and the Design of Network-based Software Architectures}, 2000.
			\item M. Fowler. \textit{Patterns of Enterprise Application Architecture}, 2002.
			\item Stripe. \textit{Checkout Sessions API Documentation}, 2024.
			\item React Team. \textit{React Documentation}, 2024.
			\item OWASP. \textit{Authentication Cheat Sheet}, 2023.
		\end{itemize}
	\end{frame}
	
\end{document}